% year: תשע"ט
% semester: ב
% moed: א
% number: 1א
% points: 10
% categories: derivatives
% source: exams/מבחנים/תשעט/חודיסמן מועד א תשעט.pdf

\begin{question}
(10 נק') חשבו את הנגזרת של הפונקציה $f(x)=\sqrt{5x+3}$ בנקודה $x$ לפי הגדרת הנגזרת (תשובה לפי טבלת הנגזרות לא תתקבל).
\end{question}

\begin{hint}
הכפילו את המונה והמכנה בצמוד $\sqrt{5(x+h)+3}+\sqrt{5x+3}$.
\end{hint}

\begin{solution}
הפונקציה מוגדרת עבור $x\ge -\tfrac35$, ונחשב את הנגזרת בנקודה $x>-\tfrac35$. לפי ההגדרה, בעזרת הכפלה בצמוד:
\begin{align*}
f'(x)&=\lim_{h\to0}\frac{\sqrt{5(x+h)+3}-\sqrt{5x+3}}{h}
=\lim_{h\to0}\frac{\big(5(x+h)+3\big)-(5x+3)}{h\left(\sqrt{5x+5h+3}+\sqrt{5x+3}\right)}\\
&=\lim_{h\to0}\frac{5}{\sqrt{5x+5h+3}+\sqrt{5x+3}}=\frac{5}{2\sqrt{5x+3}},
\end{align*}
כאשר במעבר האחרון השתמשנו ברציפות פונקציית השורש ($\sqrt{5x+5h+3}\to\sqrt{5x+3}>0$). לכן
\[ f'(x)=\frac{5}{2\sqrt{5x+3}},\qquad x>-\tfrac35. \]
(בנקודה $x=-\tfrac35$ המנה $\frac{\sqrt{5h}}{h}\to+\infty$, ולכן שם הפונקציה אינה גזירה.)
\end{solution}
